% GENERATED FILE. Edit canonical lessons, then run npm run generate:books.
% canonical-source-hash: fnv1a64:ec5d8b546203525f
% canonical-lessons: JA-C92-hoshii, JA-C92-negai, JA-C92-negau, JA-C92-tanomi, JA-C92-yokubou

\chapter{Wanted, Wish, To wish for, Request, Desire}
\label{ch:ja-wanted-wish-to-wish-for-request-desire}
\hlchaptermodality{92}

\begin{chapteropening}
\textbf{By the end of this chapter:} \emph{I can say wanted, a wish, to wish for, a request and a desire.}

Complete the last lesson of chapter 92: I can say wanted, a wish, to wish for, a request and a desire.
\end{chapteropening}

\section[\texorpdfstring{\ja{ほしい} (hoshii)}{ほしい (hoshii)}]{\ja{ほしい} (hoshii) — wanted, wanting to have}

\label{lesson:JA-C92-hoshii}



\begin{quote}
Before the new one: say the Japanese for to help, then the Japanese for to get lost.
\end{quote}

\subsection*{You'll want to know: \ja{ほしい}}
\textbf{\ja{ほしい}} — \emph{hoshii} — \textquotedblleft{}wanted, wanting to have\textquotedblright{}.

\begin{grammarlens}[title={What you've built}]
One more word for this chapter.
\end{grammarlens}

\subsection*{Your turn}
\begin{itemize}
\raggedright
  \item \emph{Say it:} \emph{hoshii}
  \item \emph{Say it:} \emph{hoshii}, once more
  \item \emph{Say it:} say \emph{hoshii}
  \item \emph{Recall:} say the Japanese for to help, then the Japanese for to get lost, then say \emph{hoshii} again
\end{itemize}

\begin{tcolorbox}[breakable,colback=teal!4,colframe=teal!35!black,title={Before you move on}]
What does \ja{ほしい} mean? (\textquotedblleft{}Wanted, wanting to have\textquotedblright{}.) Say it once more.
\end{tcolorbox}
\section[\texorpdfstring{\ja{ねがい} (negai)}{ねがい (negai)}]{\ja{ねがい} (negai) — a wish, a request}

\label{lesson:JA-C92-negai}



\begin{quote}
Before the new one: say the Japanese for to get lost, then the Japanese for wanted.
\end{quote}

\subsection*{You'll want to know: \ja{ねがい}}
\textbf{\ja{ねがい}} — \emph{negai} — \textquotedblleft{}a wish, a request\textquotedblright{}.

\begin{grammarlens}[title={What you've built}]
One more word for this chapter.
\end{grammarlens}

\subsection*{Your turn}
\begin{itemize}
\raggedright
  \item \emph{Say it:} \emph{negai}
  \item \emph{Say it:} \emph{negai}, once more
  \item \emph{Say it:} say \emph{negai}
  \item \emph{Recall:} say the Japanese for to get lost, then the Japanese for wanted, then say \emph{negai} again
\end{itemize}

\begin{tcolorbox}[breakable,colback=teal!4,colframe=teal!35!black,title={Before you move on}]
What does \ja{ねがい} mean? (\textquotedblleft{}A wish, a request\textquotedblright{}.) Say it once more.
\end{tcolorbox}
\section[\texorpdfstring{\ja{ねがう} (negau)}{ねがう (negau)}]{\ja{ねがう} (negau) — to wish for}

\label{lesson:JA-C92-negau}



\begin{quote}
Before the new one: say the Japanese for wanted, then the Japanese for a wish.
\end{quote}

\subsection*{You'll want to know: \ja{ねがう}}
\textbf{\ja{ねがう}} — \emph{negau} — \textquotedblleft{}to wish for\textquotedblright{}.

\begin{grammarlens}[title={What you've built}]
One more word for this chapter.
\end{grammarlens}

\subsection*{Your turn}
\begin{itemize}
\raggedright
  \item \emph{Say it:} \emph{negau}
  \item \emph{Say it:} \emph{negau}, once more
  \item \emph{Say it:} say \emph{negau}
  \item \emph{Recall:} say the Japanese for wanted, then the Japanese for a wish, then say \emph{negau} again
\end{itemize}

\begin{tcolorbox}[breakable,colback=teal!4,colframe=teal!35!black,title={Before you move on}]
What does \ja{ねがう} mean? (\textquotedblleft{}To wish for\textquotedblright{}.) Say it once more.
\end{tcolorbox}
\section[\texorpdfstring{\ja{たのみ} (tanomi)}{たのみ (tanomi)}]{\ja{たのみ} (tanomi) — a request}

\label{lesson:JA-C92-tanomi}



\begin{quote}
Before the new one: say the Japanese for a wish, then the Japanese for to wish for.
\end{quote}

\subsection*{You'll want to know: \ja{たのみ}}
\textbf{\ja{たのみ}} — \emph{tanomi} — \textquotedblleft{}a request\textquotedblright{}.

\begin{grammarlens}[title={What you've built}]
One more word for this chapter.
\end{grammarlens}

\subsection*{Your turn}
\begin{itemize}
\raggedright
  \item \emph{Say it:} \emph{tanomi}
  \item \emph{Say it:} \emph{tanomi}, once more
  \item \emph{Say it:} say \emph{tanomi}
  \item \emph{Recall:} say the Japanese for a wish, then the Japanese for to wish for, then say \emph{tanomi} again
\end{itemize}

\begin{tcolorbox}[breakable,colback=teal!4,colframe=teal!35!black,title={Before you move on}]
What does \ja{たのみ} mean? (\textquotedblleft{}A request\textquotedblright{}.) Say it once more.
\end{tcolorbox}
\section[\texorpdfstring{\ja{よくぼう} (yokubou)}{よくぼう (yokubou)}]{\ja{よくぼう} (yokubou) — a desire}

\label{lesson:JA-C92-yokubou}



\begin{quote}
Before the new one: say the Japanese for to wish for, then the Japanese for a request.
\end{quote}

\subsection*{You'll want to know: \ja{よくぼう}}
\textbf{\ja{よくぼう}} — \emph{yokubou} — \textquotedblleft{}a desire\textquotedblright{}.

\begin{grammarlens}[title={What you've built}]
One more word for this chapter.
\end{grammarlens}

\subsection*{Your turn}
\begin{itemize}
\raggedright
  \item \emph{Say it:} \emph{yokubou}
  \item \emph{Say it:} \emph{yokubou}, once more
  \item \emph{Say it:} say \emph{yokubou}
  \item \emph{Recall:} say the Japanese for to wish for, then the Japanese for a request, then say \emph{yokubou} again
\end{itemize}

\begin{tcolorbox}[breakable,colback=teal!4,colframe=teal!35!black,title={Before you move on}]
What does \ja{よくぼう} mean? (\textquotedblleft{}A desire\textquotedblright{}.) Say it once more.
\end{tcolorbox}
